Jedes Minispiel ist aus drei Leveln aufgebaut, welche progressiv schwerer werden.
So dient das erste Level als Einführung in die Mechaniken des jeweiligen Minispiels.
Das zweite Level kann als \glqq standard\grqq\ Schwierigkeitsgrad verstanden werden und im Dritten können die Spielenden dann ihre erlernten Fähigkeiten in schwierigen Herausforderungen unter Beweis stellen.
Nach dem erfolgreichen Beenden eines Levels erhält die Spielfigur ein Item, welches dann im Kamp verwendet werden kann.
Diese Items hängen von der Fachschaft und des Schwierigkeitsgrades ab.
Die Stärke der Items skaliert also mit der Schwierigkeit des beendeten Levels. 
Da das Bearbeiten schwererer Level länger dauert, wird dies dementsprechend belohnt.\\
Im Allgemeinen soll jedes Minispiel mit einer Auswahl des Schwierigkeitsgrades beginnen, welche auch wieder verlassen werden kann, wenn sich doch gegen das Spielen entschieden wird.
Jedes Level bietet zusätzlich noch die Möglichkeit zur Auswahl zurückzukehren.
Nach dem Abschließen des selbigen ist es ebenfalls möglich das Level neu zu starten.


\subsection{Informatik}

Die Level der Fachschaft Informatik bestehen jeweils aus einem Schaltkreis, bei dem Logikgatter fehlen und einer bestimmten Menge ebendieser, aus denen die fehlenden ausgewählt werden müssen.
Verschiedene Schwierigkeiten werden auf zwei Weisen ermöglicht:

\begin{itemize}
    \item[] Zum einen verändert sich die Anzahl an Logikgattern, aus denen ausgewählt werden kann.
            So sind es im ersten Level 4 Logikgatter, im zweiten Level 5 Logikgatter und im dritten Level 6 Logikgatter, die zur Auswahl zur Verfügung stehen.
    \item[] Zum anderen ändert sich der Schaltkreis, bei dem die Gatter entfernt wurden, sowie die Anzahl an entfernten Gattern.
\end{itemize}

Jedes Level verwendet außerdem, im Sinne der Variabilität, drei verschiedene Schaltkreise, aus denen, bei jedem Minispiel-Start, einer zufällig gewählt wird.

\subsubsection{Level 1}

Bei den Schaltkreisen des ersten Levels fehlt jeweils nur ein Logikgatter, das zwei Eingänge mit einem Ausgang verbindet.
Gesucht wird also das Logikgatter $G \in \{\land , \lor , \overline{\land}, \overline{\lor}, \veebar, \overline{\veebar}\}$, welches folgende Bedingungen erfüllt:

\begin{align*}
    x_1 G x_2 = y
\end{align*}

Dadurch, dass jeder Schaltkreis zwei Eingänge und einen Ausgang besitzt, ergeben sich vier mögliche Zustände, die in Form einer Wahrheitswerttabelle dargestellt werden können.
Für die Rätsel des ersten Levels wurden folgende Tabellen festgelegt:

\begin{table}[h!]
    \centering
    \begin{minipage}{.33\textwidth}
        $G = \land$\\
        \vfill
        \begin{tabular}{|c|c|c|} 
            \hline
            $x_1$ & $x_2$ & $y$ \\
            \hline\hline
            $0$ & $0$ & $0$ \\
            $0$ & $1$ & $0$ \\
            $1$ & $0$ & $0$ \\
            $1$ & $1$ & $1$ \\
            \hline
        \end{tabular}
    \end{minipage}
    \hfill
    \begin{minipage}{.32\textwidth}
        $G = \lor$\\
        \vfill
        \begin{tabular}{|c|c|c|} 
            \hline
            $x_1$ & $x_2$ & $y$ \\
            \hline\hline
            $0$ & $0$ & $0$ \\
            $0$ & $1$ & $1$ \\
            $1$ & $0$ & $1$ \\
            $1$ & $1$ & $1$ \\
            \hline
        \end{tabular}
    \end{minipage}
    \hfill
    \begin{minipage}{.33\textwidth}
        $G = \overline{\land}$\\
        \vfill
        \begin{tabular}{|c|c|c|} 
            \hline
            $x_1$ & $x_2$ & $y$ \\
            \hline\hline
            $0$ & $0$ & $1$ \\
            $0$ & $1$ & $1$ \\
            $1$ & $0$ & $1$ \\
            $1$ & $1$ & $0$ \\
            \hline
        \end{tabular}
    \end{minipage}
    \caption{Wahrheitswerttabellen für Level 1 des Informatik-Minispiels}
    \label{tab:wwt01}
\end{table}

Diese Tabellen werden den Spielenden, aufgrund der ansprechenderen Gestaltung, in Form einer Abbildung präsentiert, welche die vier Zustände des Schaltkreises darstellt.
Ebenfalls im Sinne der Gestaltung wurde auf die Konvention, die Eingänge auf der linken Seite und die Ausgänge auf der rechten Seite darzustellen, verzichtet.
In der Abbildung \ref{tab:wwt01} ist dies beispielhaft anhand der ersten Wahrheitswerttabelle ($G = \land$) zu erkennen.

\subsubsection{Level 2}

Im zweiten Level bestehen die Schaltkreise ebenfalls aus zwei Eingängen, sowie einem Ausgang.
Der Unterschied zu Level 1 besteht hierbei darin, dass zwei Logikgatter entfernt wurden.
Durch das erste werden beide Eingänge verbunden.
Der dadurch resultierende Ausgang wird im zweiten Logikgatter wieder mit einem der Eingänge verbunden, wodurch der finale Ausgang entsteht.
Außerdem wurde zur Reduzierung der Komplexitität die Einschränkung hinzugefügt, dass sich die Gatter unterscheiden müssen.
Gesucht werden also die Logikgatter $G_1, G_2 \in \{\land , \lor , \overline{\land}, \overline{\lor}, \veebar, \overline{\veebar}\}$, welche folgende Bedingungen erfüllen:

\begin{align*}
    (x_1 G_1 x_2) G_2 x_2 &= y\\
    G_1 &\neq G_2
\end{align*}

Dadurch, dass jeder Schaltkreis auch im zweiten Level zwei Eingänge und einen Ausgang besitzt, ergeben sich wieder vier mögliche Zustände, die in Form einer Wahrheitswerttabelle dargestellt werden können.
Für die Rätsel des zweiten Levels wurden folgende Tabellen festgelegt:

\begin{table}[h!]
    \centering
    \begin{minipage}{.33\textwidth}
        $G_1 = \land$\\
        $G_2 = \overline{\lor}$
        \vfill
        \begin{tabular}{|c|c|c|} 
            \hline
            $x_1$ & $x_2$ & $y$ \\
            \hline\hline
            $0$ & $0$ & $1$ \\
            $0$ & $1$ & $0$ \\
            $1$ & $0$ & $1$ \\
            $1$ & $1$ & $0$ \\
            \hline
        \end{tabular}
    \end{minipage}
    \hfill
    \begin{minipage}{.32\textwidth}
        $G_1 = \lor$\\
        $G_2 = \veebar$
        \vfill
        \begin{tabular}{|c|c|c|} 
            \hline
            $x_1$ & $x_2$ & $y$ \\
            \hline\hline
            $0$ & $0$ & $0$ \\
            $0$ & $1$ & $0$ \\
            $1$ & $0$ & $1$ \\
            $1$ & $1$ & $0$ \\
            \hline
        \end{tabular}
    \end{minipage}
    \hfill
    \begin{minipage}{.33\textwidth}
        $G_1 = \overline{\land}$\\
        $G_2 = \lor$
        \vfill
        \begin{tabular}{|c|c|c|} 
            \hline
            $x_1$ & $x_2$ & $y$ \\
            \hline\hline
            $0$ & $0$ & $1$ \\
            $0$ & $1$ & $1$ \\
            $1$ & $0$ & $1$ \\
            $1$ & $1$ & $1$ \\
            \hline
        \end{tabular}
    \end{minipage}
    \caption{Wahrheitswerttabellen für Level 2 des Informatik-Minispiels}
    \label{tab:wwt02}
\end{table}

Diese Tabellen werden analog zu Level 1 als Abbildungen dargestellt.

\subsubsection{Level 3}

Im dritten Level bestehen die Schaltkreise aus drei Eingängen, einem Ausgang und drei entfernten Logikgattern.
Durch die ersten beiden werden jeweils zwei der Eingänge miteinander verbunden.
Die dadurch resultierenden Ausgänge werden im dritten Logikgatter verbunden, wodurch der finale Ausgang entsteht.
Gesucht werden also die Logikgatter $G_1, G_2, G_3 \in \{\land , \lor , \overline{\land}, \overline{\lor}, \veebar, \overline{\veebar}\}$, welche folgende Bedingungen erfüllen:

\begin{align*}
    (x_1 G_1 x_2) G_3 (x_2 G_2 x_3) &= y\\
    G_1 \neq G_2 &\neq G_3 
\end{align*}

Dadurch, dass jeder Schaltkreis aus drei Eingänge und einen Ausgang besteht, ergeben sich insgesamt acht mögliche Zustände, die in Form einer Wahrheitswerttabelle dargestellt werden können.
Für die Rätsel des dritten Levels wurden folgende Tabellen festgelegt:

\begin{table}[h!]
    \centering
    \begin{minipage}{.33\textwidth}
        $G_1 = \land$\\
        $G_2 = \veebar$\\
        $G_3 = \lor$
        \vfill
        \begin{tabular}{|c|c|c|c|} 
            \hline
            $x_1$ & $x_2$ & $x_3$ & $y$\\
            \hline\hline
            $0$ & $0$ & $0$ & $0$\\
            $0$ & $0$ & $1$ & $1$\\
            $0$ & $1$ & $0$ & $1$\\
            $0$ & $1$ & $1$ & $0$\\
            $1$ & $0$ & $0$ & $0$\\
            $1$ & $0$ & $1$ & $1$\\
            $1$ & $1$ & $0$ & $1$\\
            $1$ & $1$ & $1$ & $1$\\
            \hline
        \end{tabular}
    \end{minipage}
    \hfill
    \begin{minipage}{.32\textwidth}
        $G_1 = \overline{\veebar}$\\
        $G_2 = \land$\\
        $G_3 = \lor$
        \vfill
        \begin{tabular}{|c|c|c|c|} 
            \hline
            $x_1$ & $x_2$ & $x_3$ & $y$\\
            \hline\hline
            $0$ & $0$ & $0$ & $1$\\
            $0$ & $0$ & $1$ & $1$\\
            $0$ & $1$ & $0$ & $0$\\
            $0$ & $1$ & $1$ & $1$\\
            $1$ & $0$ & $0$ & $0$\\
            $1$ & $0$ & $1$ & $0$\\
            $1$ & $1$ & $0$ & $1$\\
            $1$ & $1$ & $1$ & $1$\\
            \hline
        \end{tabular}
    \end{minipage}
    \hfill
    \begin{minipage}{.33\textwidth}
        $G_1 = \veebar$\\
        $G_2 = \lor$\\
        $G_3 = \land$
        \vfill
        \begin{tabular}{|c|c|c|c|} 
            \hline
            $x_1$ & $x_2$ & $x_3$ & $y$\\
            \hline\hline
            $0$ & $0$ & $0$ & $0$\\
            $0$ & $0$ & $1$ & $0$\\
            $0$ & $1$ & $0$ & $1$\\
            $0$ & $1$ & $1$ & $1$\\
            $1$ & $0$ & $0$ & $0$\\
            $1$ & $0$ & $1$ & $1$\\
            $1$ & $1$ & $0$ & $0$\\
            $1$ & $1$ & $1$ & $0$\\
            \hline
        \end{tabular}
    \end{minipage}
    \caption{Wahrheitswerttabellen für Level 3 des Informatik-Minispiels}
    \label{tab:wwt03}
\end{table}

Da die Wahrheitswerttabellen des dritten Levels acht Zustände beinhalten, wird auf die Darstellung jedes einzelnen Zustandes verzichtet und nur der erste Zustand als Schaltkreis abgebildet.
Die anderen werden als Tabelle präsentiert.